\begin{statement}{A.13}{Proved}{Punctured elliptic maps for the quaternion and stable dihedral covers}
Work over an algebraic closure of $\mathbf Q(i)$; every curve, map,
origin, and isogeny below is defined over a finite extension.
Let $S$ be the cusp set specified in A.4.

For the quaternion cover
\[
 C_Q:X^4+Y^4=2Z^4,\qquad S=\{Z=0\},
\]
each of the three sign-involution elliptic quotients can be followed
by an isogeny so that the inverse image of its target origin is
\emph{exactly} $S$. For the sign changes of $X$ and $Y$ these
composites have degree four and have target
$E':B^2=A^3-32A$, with affine formulas
\[
 \phi_X(x,y)=(4y^2,-8i\,yx^2),\qquad
 \phi_Y(x,y)=(4x^2,-8i\,xy^2),\qquad x=X/Z,\quad y=Y/Z.
                                                        \tag{A.13.1}
\]
Their fibres over $O$ are reduced. For the sign change of $Z$,
choose one cusp image as origin on the elliptic quotient and
follow its degree-two quotient map by multiplication by two.
The resulting degree-eight map has fibre $2\sum_{P\in S}P$
over the origin. All three elliptic targets have CM by
$\mathbf Z[i]$, and their pulled-back differential lines span
$H^0(C_Q,\Omega^1)$.

For the stable dihedral genus-two curve
\[
 D_c:z^2=(t^2-4c)(t^4-4ct^2+2c^2),\qquad c^4=-1,
\]
write $Q_+,Q_-$ for its two points at infinity. There is a
degree-two map
\[
 D_c\longrightarrow E_c^{-2}:\quad (t,z)\longmapsto(X,Y)=(t^2,z),
 \qquad
 Y^2=X^3-8cX^2+18c^2X-8c^3.                         \tag{A.13.2}
\]
This elliptic curve has $j=8000$ and CM by
$\mathbf Z[\sqrt{-2}]$. The origin fibre is precisely
$Q_++Q_-$. It is the complementary elliptic factor to the
$\mathbf Q(i)$ factor of A.4 and defines a punctured map with
no extra points removed from the source.

For the stable symmetric-group genus-two curve
\[
 D_3:z^2=(t^2+4)(t^3+3t),\qquad P=\infty,
\]
every nonconstant morphism to an elliptic curve whose origin fibre
is supported only at $P$ has degree three and ramification divisor
$2P$. All such maps pull back elliptic differentials to the same
one-dimensional subspace of $H^0(D_3,\Omega^1)$; they cannot detect
the second elliptic factor independently. In fact they coincide
up to an isomorphism of their pointed elliptic targets.
Thus the proposed second independent punctured elliptic map
for this one-cusp genus-two curve does not exist.
\end{statement}
\begin{proof}
For the sign change of $X$, the quotient has coordinates
$t=Y/Z$, $u=X^2/Z^2$ and equation
\[
                         u^2=2-t^4.
\]
This is a smooth genus-one quartic. Its two points at infinity,
$P_+$ and $P_-$, have $u/t^2=+i$ and $-i$, respectively.
The sign involution has no fixed cusp, so the four cusps map
in pairs to these two points and are unramified there.
The identity
\[
                   (u-it^2)(u+it^2)=2
\]
and the expansions at infinity give
\[
             \operatorname{div}(u-it^2)=2P_+-2P_-.
\]
The difference $P_--P_+$ is nonzero: otherwise a function with
one simple pole would give a degree-one map from this elliptic
curve to a projective line. Choose $P_+$ as origin. The other
point is therefore an exact two-torsion point, and the isogeny
with kernel these two points has precisely their union as its
origin fibre. Pulling back gives exactly the four original cusps.

There are explicit coordinates for this isogeny. Put
$v=-iu$, so $v^2=t^4-2$, and set
\[
       a=2(v+t^2),\qquad b=4t(v+t^2).
\]
Then $b^2=a^3+8a$, the chosen origin is infinity, and the other
infinity is $(0,0)$. The standard degree-two isogeny with kernel
$(0,0)$ is
\[
              (a,b)\longmapsto
                   \left(a+\frac8a,b\left(1-\frac8{a^2}\right)\right).
\]
Its two coordinates simplify to $4t^2$ and $8tv$.
This proves the first formula in (A.13.1); interchange $X,Y$ to
obtain the second. The target has the automorphism
$(A,B)\mapsto(-A,iB)$, and is isomorphic to $y^2=x^3-x$
after adjoining a fourth root of $32$.

For the sign change of $Z$, use $t=X/Y$, $u=Z^2/Y^2$ on
the quotient. Its equation is $2u^2=t^4+1$. The four cusp
images are exactly the four points $u=0$. Choose one of these
as origin. The involution $u\mapsto-u$ fixes this origin and
is the elliptic involution $[-1]$; hence those four points are
exactly the full two-torsion subgroup. Multiplication by two
has precisely that subgroup as its origin fibre. Each original
cusp is ramified with index two in the sign-$Z$ quotient, so the
composite fibre has multiplicity two at each of the four cusps.
The quartic has $j=1728$, either by its order-four automorphism
or by the elementary quartic invariant formula. Its target under
$[2]$ is the same elliptic curve.

Finally the canonical forms on the plane quartic are the
adjunction forms corresponding to $X,Y,Z$. Under the three
sign involutions their respective invariant lines are distinct
and together span the three-dimensional differential space.
These are the pullbacks from the three genus-one quotients.
An isogeny in characteristic zero has nonzero differential,
so composing with the displayed isogenies leaves those three
lines unchanged. This proves the claimed simultaneous detection
of all three proper elliptic factors.

For (A.13.2), direct substitution $X=t^2$, $Y=z$ gives the cubic.
Its discriminant and invariant are
\[
                       \Delta=512c^6,\qquad c_4=160c^2,
\]
so it is smooth and $j=c_4^3/\Delta=8000$. At each infinity
of $D_c$, the functions $X$ and $Y$ have pole orders two and
three, respectively. Thus each point maps unramified to the
cubic's origin, and their total degree two exhausts the fibre.

The CM assertion can also be proved without consulting a table
of singular moduli. Translate $x=X-4c$ to write the cubic as
\[
                         y^2=x^3+4cx^2+2c^2x.
\]
The two-isogeny with kernel $(0,0)$ has target
$y'^2=x'^3-8cx'^2+8c^2x'$. This target is isomorphic to the
source by $(x',y')=(-2x,2\sqrt{-2}\,y)$.
Compose the two-isogeny with the inverse isomorphism. The resulting
degree-two endomorphism has differential $-\sqrt{-2}$, so its
square is multiplication by $-2$: endomorphisms of an elliptic
curve in characteristic zero are determined by their differentials.
This embeds the maximal quadratic order $\mathbf Z[\sqrt{-2}]$
in its endomorphism ring. It is therefore the CM ring.
It differs from $\mathbf Q(i)$, so this elliptic quotient is
the complementary factor in A.4.

For the last assertion, the point at infinity on the odd-degree
hyperelliptic model $D_3$ is a Weierstrass point. If
$f:D_3\to E_0$ is nonconstant and $f^{-1}(O)$ is supported
only there, then $f^*O=dP$, where $d=\deg f$. The ramification
divisor $R_f$ has degree two by Riemann--Hurwitz and contains
$(d-1)P$. Thus $2\le d\le3$. The case $d=2$ is impossible:
$R_f$ is the divisor of the pullback of a nonzero elliptic
differential, hence is canonical. The unique canonical divisor
containing the Weierstrass point $P$ is $2P$, whereas a
degree-two map has ramification multiplicity at most one at
each point. Therefore $d=3$ and $R_f=2P$.
Consequently
\[
                 f^*H^0(E_0,\Omega^1)=H^0(D_3,\Omega^1(-2P)),
\]
a one-dimensional space. The degree-three map already supplied
by A.4 gives precisely this line.

For completeness, equality of these lines gives uniqueness of
the pointed target quotient. Embed the field of definition in
$\mathbf C$. The induced homomorphism
$J(D_3)\to E_0$ has connected kernel. Otherwise $f$ factors
through a nontrivial finite étale isogeny of elliptic curves;
its degree divides three, and the remaining map would have
degree one, which a genus-two curve cannot have. The connected
kernel is the elliptic subvariety whose tangent space is the
annihilator of the displayed differential line. A connected
complex subtorus is determined by its tangent space. Hence all
these kernels are equal, and the quotient elliptic curves are
isomorphic. Since $P$ maps to the chosen origins, no translation
ambiguity remains.

All displayed functions extend to morphisms of smooth proper
curves, and the exact fibre computations show that they restrict
to morphisms of the indicated punctured curves. After compatible
geometric base points and paths are chosen, they induce the
usual continuous homomorphisms of arithmetic and geometric
fundamental groups. A change of path conjugates these maps by
a geometric element; the proper-homology detection assertions
are unaffected. The executable certificate
\texttt{code/a\_cover\_maps.py} verifies the map identities,
the two-isogeny identities, and the elliptic invariants.
\end{proof}
